\chapter{Fondamenti di algebra lineare}
L'algebra lineare è il linguaggio fondamentale della meccanica quantistica, in cui gli stati fisici sono rappresentati come vettori in uno spazio di Hilbert complesso. Operatori hermitiani agiscono su questi vettori per rappresentare osservabili fisiche, mentre le probabilità di misura sono calcolate tramite prodotti scalari, rendendo la teoria quasi interamente formulabile in termini di algebra lineare. Introduciamo quindi i fondamenti di algebra linearre che saranno necessari per la comprensione del corso.



\section{Spazio di Hilbert e prodotto scalare}
Si dice \textbf{spazio di Hilbert} uno spazio vettoriale completo rispetto alla norma indotta dal prodotto scalare. Ci saranno utili nello studio dei sistemi quantistici, a causa dell'assioma 1 della meccanica quantistica.

\begin{center}
	\vspace{3pt}
	\textbf{\Large Assioma 1 della meccanica quantistica}\\
	\vspace{5pt}
	\it Ad ogni sistema fisico quantistico è associato uno spazio di Hilbert $\mathcal{H}$.\rm
\end{center}

Tutti gli spazi di Hilbert che andremo a trattare sono finiti, in quanto i sistemi su cui operiamo sono finiti.
Consideriamo lo \textbf{spazio di Hilbert} $\mathbb{C}^d$ dei vettori colonna sui numeri complessi. Un elemento di questo spazio di Hilbert, è del tipo

$$
\overline{v} =
\begin{pmatrix}
	v_1 \\
	v_2 \\
	\vdots \\
	v_d
\end{pmatrix},
\hspace{0.3cm} v_i \in \mathbb{C}
$$
È inoltre dotato del \textbf{prodotto scalare} $\braket{v|w} = \sum^{d-1}_{i=0}\bar{v_i}w_i$.
Nel contesto dei numeri complessi, $\overline v$ indica il coniugato complesso di $v$, ossia:
$$
v = a + ib, \qquad \overline v = a - ib
$$
%
Sia $\mathcal{H}$ è uno spazio di Hilbert arbitrario con $\dim(\mathcal{H} = d)$ finita, posso scegliere una \textbf{base ortonormale} $\left\{ \ket{e_i} \right\}^d_{i=1}$ di $\mathcal{H}$, ossia tale che:

$$
\braket{e_i|e_j}  = 
\begin{cases}
	1 & \text{se }i=j  \\
	0 & \text{se }i\neq j
\end{cases}
$$
%
Indichiamo i vettori di $\mathcal{H}$ con il simbolo $\ket{\psi}$ e $\bra{\psi}$ per il suo duale. $\ket{\psi}$ è un vettore colonna, mentre $\bra{\psi}$ è un vettore riga. 

\paragraph{Base ortonormale standard} Una base ortonormale particolare, è la \textbf{base ortonormale standard}.
La base ortonormale standard di un qualsiasi spazio di Hilbert $\mathcal{H}= \mathbb{C}^d$ è data dall'insieme di vettori
$$
\ket{0}=
\begin{pmatrix}
	1 \\
	0 \\
	\vdots \\
	0
\end{pmatrix}, \
\ket{1}=
\begin{pmatrix}
	0 \\
	1 \\
	\vdots \\
	0
\end{pmatrix}, \dots,\
\ket{d-1}=
\begin{pmatrix}
	0 \\
	0 \\
	\vdots \\
	1
\end{pmatrix}
$$

Se $\ket{\psi} = 
\begin{pmatrix}
	\psi_1 \\
	\psi_2 \\
	\vdots \\
	\psi_{d-1}
\end{pmatrix}
$, allora $\ket{\psi} = \sum^{d-1}_{i=0}\psi_i \ket{i}$, e il suo duale è:
$\bra{\psi} = \left(\overline{\psi}_0,\overline{\psi}_1,\dots,\overline{\psi}_{d-1} \right)$. 

\subsection{Disuguaglianza di Schwartz}
Il prodotto scalare soddisfa la seguente disuguaglianza, nota come disuguaglianza
$$\forall \ket{\Phi},\ket{\Psi} \in \mathcal{H} = \mathbb{C}^d, \hspace{0.2cm} \left|\braket{\Phi|\Psi} \right|^2 \le \braket{\Phi|\Phi} \cdot \braket{\Psi|\Psi}$$
E vale solo se: $$\ket{\Phi} = k\ket{\Psi}, \exists k \in\mathbb{C}$$

La norma indotta del prodotto scalare è $\left| \left| \Psi\right| \right| = \sqrt{\braket{\Psi|\Psi}}$, per cui possiamo riscrivere il lemma di Schwartz come:
$$\braket{\Phi|\Psi} = \left| \left|\Phi \right| \right| \cdot \left| \left|\Psi \right| \right|$$

\section{Mappe lineari e matrici}
Utilizziamo funzioni lineari $f$, con le seguenti proprietà:
\begin{enumerate}
	\item $f(\lambda x) = \lambda f(x)$,
	\item $f(x+y) = f(x) + f(y)$ 
\end{enumerate}

Dati $\mathcal{H}= \mathbb{C}^d, \mathcal{K} = \mathbb{C}^e$ spazi di Hilbert, $\textit{lin}(\mathcal{H},\mathcal{K})$ è l'insieme delle \textbf{mappe lineari} da $\mathcal{H}$ a $\mathcal{K}$.\\ Se fissiamo una base per $\mathcal{H}$ e una base per $\mathcal{K}$ allora ogni mappa di $\textit{lin}(\mathcal{H},\mathcal{K})$ può essere scritta come una matrice $d\times e$.

Quando la mappa lineare va dallo stesso spazio allo stesso spazio, si dice \textbf{endomorfismo}. Un esempio è $\textit{lin}(\mathcal{H},\mathcal{H}) $, che indicheremo più brevemente con $\textit{lin}(\mathcal{H})$. 


Una particolare mappa lineare $\mathcal{H} \to \mathcal{H}$ è l'identità $ \mathbb{1}$.

\subsection{Traccia}
La \textbf{Traccia} di $M \in \textit{lin}(\mathcal{H})$ si calcola scegliendo una base $\Sigma$ di $\mathcal{H}$ e sommando gli elementi della diagonale principale.
$$\text{tr}[M] = \sum_{i \in \Sigma} \braket{i|M|i}$$
La traccia  è un invariante  per cambiamenti di base e soddisfa 
$$\text{tr}[MN] = \text{tr}[NM],\hspace{0.4cm} \forall M\in \textit{lin}(\mathcal{H,K}),N\in \textit{lin}(\mathcal{K,H})$$

Se $M\in \textit{lin}(\mathcal{H})$ e $N = \ket{\Psi}\bra{\Psi}$ per $\ket{\Psi}\in \mathcal{H}$ allora 
$\text{tr}[MN]=\braket{\Psi|M|\Psi}$

\subsection{Autovettori e autovalori}
Se $M\in \textit{lin}(\mathcal{H})$ e $\exists\ket{\Psi}\in \mathcal{H}$ non nullo e $n \in \mathbb{C}$ t.c $M \ket{\Psi} = \lambda \ket{\Psi}$ allora $\ket{\Psi}$ è un \textbf{Autovettore} di $M$ con \textbf{Autovalore} $\lambda$.

Se $M\in \textit{lin}(\mathcal{H}_1,\mathcal{H}_2)$, dove $\mathcal{H}_1$ ha base $\Sigma_1$ e  $\mathcal{H}_2$ ha base $\Sigma_2$ allora 
$$M = \sum_{i \in \Sigma_1}\sum_{j \in \Sigma_2} M_{ij} \ket{i}\bra{j}\hspace{0.2cm} \text{e} \hspace{0.2cm} M_{ij} = \braket{i|M|j}$$

\subsection{Matrice trasposta coniugata (o aggiunta, o dagger)}
Sia  $M\in \textit{lin}(\mathcal{H}_1,\mathcal{H}_2)$, l'\textbf{operatore aggiunto} di $M$ è $M^\dagger$ (dagger), tale che $M^\dagger = \overline{M}^T$, e che

$$
\braket{\Psi|M|\Phi} = \overline{\braket{\Phi|M^\dagger|\Psi}}
$$

Ed inoltre
$$M^\dagger\in \textit{lin}(\mathcal{H}_2,\mathcal{H}_1) \qquad \text{e} \qquad M^\dagger = \sum_{i,j} \overline{M_{ij}} \cdot\ket{j}\bra{i}$$

Una proprietà interessante, è che, $\forall M \in \textit{lin}(\mathcal H_2, \mathcal H_3),\; N \in \textit{lin}(\mathcal H_1, \mathcal H_2)$, vale che $(MN)^\dagger = N^\dagger M^\dagger$.

\subsection{Operatori Hermitiani}
$M\in \textit{lin}(\mathcal{H})$ è \textbf{Hermitiano} (o autoaggiunto) se $M = M^\dagger$, cioè significa che $M$ ha elementi diagonali $M_{ii}$ reali e $M_{ij} = \overline{M_{ij}}$ per $i\neq j$.

$$\begin{pmatrix}
	1 & 1+i \\
	1-i & 0\\
\end{pmatrix}$$
Le matrici Hermitiane hanno autovalori reali e ammettono una base di autovettori.

\subsection{Teorema Spettrale}
Sia $M\in \textit{lin}(\mathcal{H})$ Hermitiana e $\dim (\mathcal{H} = d)$, allora 
$\exists \lambda_1\ge\lambda_2\ge\dots\ge \lambda_d \in \mathbb{R}$ e una base $\left\{\ket{\Psi}\right\} ^d_{i=1} $ t.c.
$$M = \sum^d_{i=1}\lambda_i \ket{\Psi_i}\bra{\Psi_i}\hspace{0.4cm} \text{e\hspace{0.4cm} $\left\{ \lambda\right\} ^d_{i=1} $ si chiama \textbf{Spettro di M}}$$

Questo teorema ci dice che una matrice hermitiana può essere diagonalizzata usando matrici unitarie e la matrice diagonale ha elementi reali.
\medskip
\\OSS: lo spettro è univocamente è univocamente determinato da $M$ ma gli autovalori potrebbero essere non unici se lo spettro è degenere.
\\esempio: $\mathbb{1} = \sum^d_{i=1}\ket{\Psi_i}\bra{\Psi_i}, \hspace{0.2cm} \forall \left\{ \ket{\Psi_i}\right\} ^d_{i=1}$ 

\medskip
Se $M$ è Hermitiano  (o unitario) e $f$ è una funzione a una variabile, allora
$$f(M) = \sum^d_{i=1} f(\lambda_i)\ket{\Psi_i} \bra{\Psi_i}$$
Ossia: bisogna che $f$ sia ben definita sullo spettro.

\subsection{Operatori Positivi}
$P \in \textit{lin}(\mathcal{H})$ è \textbf{Semidefinito positivo} (\textit{PSD}) oppure solo \textit{positivo} se
$$\forall \ket{\Psi} \in \mathcal{H}, \braket{\Psi|P|\Psi} \ge 0$$
\textbf{PSD($\mathcal{H})$} è l'insieme degli operatori $P\ge 0$ su $\mathcal{H}$
\medskip

$P \in \textit{lin}(\mathcal{H})$ è \textbf{Definitivo positivo} (\textit{PD}) $P > 0$ se $\forall \braket{\Psi|P| \Psi} > 0$

\medskip
\textit{PD($\mathcal{H}$)} è l'insieme degli operatori $P > 0$ su $\mathcal{H}$

\medskip 
\textbf{Lemma (caratterizzazione degli operatori PSD)}:

$P \in \textit{lin}(\mathcal{H})$ equivale a dire:
\begin{itemize}
	\item[\textbf{A.}] $P$ è \textit{PSD}, $\braket{\Psi|P|\Psi} \ge 0 \hspace{0.5cm} \forall\ket{\Psi}\in\mathcal{H}$ 
	\item[\textbf{B.}] $P$ è Hermitiano e i suoi autovalori sono \textbf{non negativi}
	\item[\textbf{C.}] $\exists M \in \textit{lin}(\mathcal{H,K}):P=M^\dagger M \hspace{0.3cm} \exists \mathcal{K}$ spazio di Hilbert
	\item[\textbf{D.}] $\forall Q \in PSD(\mathcal{H}),\  \text{tr}[PQ] \ge 0$
\end{itemize}

\subsubsection{Corollario}
Se $$P \in PSD(\mathcal{H}), M \in \textit{lin}(\mathcal{H,K})$$ allora $$MPM^\dagger \in PSD(\mathcal{K}) \qquad P\le Q \Leftrightarrow Q-P \ge 0$$
Questa nozione induce un ordinamento...

\subsubsection{Matrici Unitarie e Isometrie}
\begin{itemize}
	\item $V \in \textit{lin}(\mathcal{H,K})$ è \textbf{Isometria} se $V^\dagger V = \mathbb{1}$ (solo se $\dim(\mathcal{H}) \le\dim(\mathcal{K})$)
	
	\item $V\in \textit{lin}(\mathcal{H})$ è \textbf{Unitaria} se $U^\dagger U = UU^\dagger = \mathbb{1}$
	
\end{itemize}

Definiamo $U(\mathcal{H})$ l'\textbf{insieme delle matrici unitarie} su $\mathcal{H}$:
$$\textit{Isom}(\mathcal{H,K})$$ 
insieme delle matrici isometrie da $\mathcal{H}$ a $\mathcal{K}$

$P \in \textit{lin}(\mathcal{H})$ è una \textbf{Proiezione} se $P^2 = P$ e $P$ è Hermitiana.

Sia $$\left\{\ket{e_i} \right\}^d_{i=1}$$ base per l'immagine di una proiezione $P$, allora 
$$P = \sum^d_{i=1}\ket{e_i}\bra{e_i}$$

Supponiamo $M$ Hermitiana, allora $M = \sum^d_{i=1} \lambda_i \ket{\Psi_i}\bra{\Psi_i}$

\medskip
I proiettori sono $P_\lambda = \sum_{i=\lambda} \ket{\Psi_i}\bra{\Psi_i}$

\medskip
Se $\Lambda =$ insieme degli autovalori \underline{distinti}, allora $$M = \sum_{\lambda\in \Lambda} \lambda P_\lambda$$

In generale $$M\in\textit{lin}(\mathcal{H,K}) \quad \text{,} \quad M = \sum^r_{i=1} s_i\ket{e_i}\bra{f_i}$$
con $s_i\ge0$ autovalori, $r = \textit{rank}(M)$ ed $e_i,f_i$ basi di $\mathcal{H}$ e $\mathcal{K}$.

\newpage
\subsection{Prodotti tensoriali}
Def. $\mathcal{H,K}$ spazi di Hilbert, il loro \textbf{Prodotto tensoriale $\mathcal{H} \otimes \mathcal{K}$} è lo spazio di Hilbert che contiene tutte le combinazioni lineari di vettori della forma $v \otimes u, \hspace{0.2cm} v\in\mathcal{H}, w \in \mathcal{K}$
\begin{itemize}
	\item $(\alpha v) \otimes (\beta w) = \alpha\beta \cdot v\otimes w, \hspace{0.2cm} \forall\alpha,\beta \in \mathbb{C}$
	\item $(v_1 + v_2) \otimes w = v_1 \otimes w + v_2 \otimes w$
	\item $v + (w_1 \otimes w_2) = v \otimes w_1 + v \otimes w_2$
\end{itemize}
Il prodotto scalare è definito da
$$\braket{v_1 \otimes w_2 |v_2 \otimes w_2} = \braket{v_1|v_2} \otimes\braket{w_1|w_2}$$
Base di $\mathcal{H} \otimes \mathcal{K} \hspace{0.5cm} \left\{ \ket{i} \otimes\ket{j} \right\}, \ket{i}\in B_\mathcal{H}, \ket{j} \in B_\mathcal{K}$

\medskip
$\ket{\Phi}\ket{\Psi} = \ket{\Phi\Psi} = \ket{\Phi} \otimes \ket{\Psi}$

\medskip
$\ket{\Phi} ^{\otimes n} = \ket{\Phi} \otimes \dots \otimes \ket{\Phi} \hspace{0.5cm} \mathcal{H}^{\otimes n} = \mathcal{H} \otimes\dots\otimes\mathcal{H}$

\medskip
Se $M \in \textit{lin}(\mathcal{H}_1,\mathcal{H}_2)$ e $N \in \textit{lin}(\mathcal{K}_1,\mathcal{K}_2)$ allora
$$M\otimes N : \mathcal{H}_1 \otimes \mathcal{K}_1 \longrightarrow \mathcal{H}_2 \otimes \mathcal{K}_2 $$
$$(M \otimes N) (\ket{ \Phi} \otimes \ket{\Psi}) = M\ket{\Phi} \otimes N \ket{\Psi}$$
$$\textit{lin}(\mathcal{H}_1\otimes\mathcal{K}_1,\mathcal{H}_2\otimes\mathcal{K}_2) = \textit{lin}(\mathcal{H}_1,\mathcal{H}_2) \otimes \textit{lin}(\mathcal{K}_1,\mathcal{K}_2)$$
$$M = \sum_{i,j}M_{ij}\ket{i}\bra{j}, \quad N = \sum_{k,l} N_{kl} \ket{k}\bra{l}$$
$$M \otimes N = \sum_{i,j,k,l} M_{ij} N_{kl}\ket{i}\bra{j} \otimes \ket{k}\bra{l} = \sum_{i,j,k,l} M_{ij} N_{kl}\ket{ik}\bra{jl}$$

\textbf{Lemma 4.1}
\begin{description}
	\item[a.] $P \in \textit{PSD}(\mathcal{H}), Q \in \textit{PSD}(\mathcal{K}) \Rightarrow P \otimes Q \in \textit{PSD}(\mathcal{H}\otimes\mathcal{K})$
	\item[b.] $\forall M \in \textit{lin}(\mathcal{H}), N \in \textit{lin}(\mathcal{K}), \text{ tr}[M\otimes N] = \text{tr}[M]\text{tr}[N]$
	\item[c.] $\forall M \in \textit{lin}(\mathcal{H}), N \in \textit{lin}(\mathcal{K}), \text{ rank}(M\otimes N) = \text{rank}(M)\text{rank}(N)$
	\item[d.] $\ket{\Phi}\otimes z = z\ket{\Phi}, z \in \mathbb{C} \hspace{1cm}\mathcal{H} \otimes \mathbb{C}= \mathcal{H}$
	\item[e.] $M\otimes \ket{\Psi}: \mathcal{H} \longrightarrow \mathcal{H} \otimes \mathcal{K} \hspace{0.8cm} (M\otimes \ket{\Psi})\ket{\Phi} = M \ket{\Phi} \otimes \ket{\Psi}$
	
	\medskip
	$M\otimes \bra{\Psi}: \mathcal{H} \otimes \mathcal{K} \longrightarrow \mathcal{H} \hspace{0.8cm} (M\otimes \bra{\Psi}) (\ket{\Phi}\otimes\bra{\chi}) = \braket{\Psi|\chi} M \ket{\Phi} \otimes \in \mathcal{H}$
\end{description}

